\chapitre{Alignement par fronts d’onde sur accélérateur}
L'implémentation eWFA-GPU de 2022~\cite{wfagpu} distribue les paires entre blocs et avance les diagonales d'un même front en parallèle. Son format de trace réserve deux bits à chacune des 64 opérations d'édition possibles. La boucle de cette version s'arrête avant le score 64. L'expérience retient les cas \(d<64\) parmi vingt paires définies avant le calcul. Pour chaque chemin brut, le programme rejoue les opérations jusqu'aux deux extrémités et vérifie
\[
 \#\text{substitutions}+\#\text{insertions}+\#\text{délétions}=d_{\rm edlib}.
\]
\chart{20_front_gpu.png}{Sur les paires admissibles, les temps de processus comparent trois méthodes avec reconstruction du chemin sur les mêmes lots.}{fig:ewfa}{bench/ewfa.sbatch}
La figure~\ref{fig:ewfa} porte sur cinq processus par lot. Pour cinq paires de 5\,000 bases faiblement divergentes, le processus eWFA-GPU prend 297,29\,ms, dont 2,93\,ms dans son chronomètre interne. NvPD CUDA prend 376,73\,ms et edlib sur un cœur Grace 6,03\,ms, du lancement du processus à sa fin. Sur les quatorze paires admissibles, le score concorde avec edlib et le rejeu du chemin brut consomme les deux chaînes à ce coût. Le formateur de sortie du programme amont produit quatorze CIGAR qui échouent au rejeu.
\synthese{Les cinq paires entières de 5\,000 bases à 2\,\% donnent des distances de 97 à 99, au-delà de la portée du format de trace.}{Pour les charges denses de plus grande distance, la récurrence par lignes se porte aussi entre architectures.}
